\documentclass[12pt]{article}
\usepackage[utf8]{inputenc}
\usepackage[T1]{fontenc}
\usepackage[spanish]{babel}
\usepackage{geometry}
\usepackage{amsmath,amssymb}
\usepackage{enumitem}

\geometry{margin=2.5cm}

\newcommand{\dudoso}[1]{\textsf{#1}}

\begin{document}

\begin{center}
{\Large Transcripción cruda}\\[2mm]
{\large 12May26}
\end{center}

\section*{Foto 1}

\textit{Ejemplo.} Sean $a,b\in\mathbb R$, $a<b$.

\[
C[a,b]=\{f:[a,b]\to\mathbb R \mid f \text{ es continua}\}
\]

Sea
\[
d:C[a,b]\times C[a,b]\to [0,\infty)
\]
\[
d(f,g)=\sup\{\,|f(x)-g(x)| \mid x\in[a,b]\,\}.
\]

Afirmación: $d$ es métrica en $C[a,b]$.

\textit{Dem.}
\begin{enumerate}[label=\arabic*)]
    \item Positividad.
    \item Simetría.
    \item Desigualdad del triángulo.
\end{enumerate}

\section*{Foto 2}

Si $f,g\in C[a,b]$, $f\neq g$, entonces $\exists x_0\in[a,b]$ tal que
\[
f(x_0)\neq g(x_0).
\]
Como $f$ y $g$ son continuas, existe un intervalo $U$ en $[a,b]$ tal que
\[
x_0\in U
\qquad \text{y}\qquad
f(x)\neq g(x),\ \forall x\in U.
\]
Entonces
\[
d(f,g)=\sup\{\,|f(x)-g(x)|\mid x\in[a,b]\,\}
\]
\[
\ge \sup\{\,|f(x)-g(x)|\mid x\in U\,\}
\]
\[
\ge |f(x_0)-g(x_0)|>0.
\]

Si $f=g$ entonces
\[
d(f,g)=d(f,f)=\sup\{\,|f(x)-g(x)|\mid x\in[a,b]\,\}=0.
\]

\section*{Foto 3}

\[
|f(x)-g(x)|=|g(x)-f(x)|\qquad \forall x\in[a,b].
\]
\[
\therefore\ d(f,g)=d(g,f).
\]

Observemos que, por la desigualdad del triángulo para $\mathbb R$,
\[
|f(x)-g(x)|\le |f(x)-h(x)|+|h(x)-g(x)|
\qquad \forall x\in[a,b].
\]
Tomando supremo sobre $x\in[a,b]$:
\[
d(f,g)\le \sup\{|f(x)-g(x)|:x\in[a,b]\}.
\]

\section*{Foto 4}

\[
d(f,g)\le \sup\{|f(x)-h(x)|:x\in[a,b]\}+\sup\{|h(x)-g(x)|:x\in[a,b]\}
\]
\[
=d(f,h)+d(h,g).
\]

\textit{Ejercicio:} Encontrar un ejemplo en donde la desigualdad señalada con $?$ es estricta.

\section*{Foto 5}

\[
d(f,g)\le d(f,h)+d(h,g).
\]

\textit{Ejercicio:} Encontrar un ejemplo en donde la desigualdad señalada con $?$ es estricta.

\section*{Foto 6}

\[
(C[a,b],d)\text{ es espacio métrico.}
\]

\textit{Ejemplo.} Sea $A$ un conjunto cualquiera. Sea $(X,d)$ un espacio métrico.
Sea
\[
\mathcal B(A,X)=\{f:A\to X\mid f \text{ es función acotada}\}.
\]
Sea
\[
D:\mathcal B(A,X)\times \mathcal B(A,X)\to [0,\infty)
\]
\[
D(f,g)=\sup\{d(f(x),g(x))\mid x\in A\}.
\]
Donde por función acotada, lo siguiente:

\section*{Foto 7}

Si $g\in X$ y $f:A\to X$ es una función, dada por
\[
\gamma(a)=d(f(a),g),
\]
entonces $\gamma$ es acotada.

Observemos lo siguiente:

\[
\gamma:A\to\mathbb R,
\qquad
\gamma(x)=d(f(x),g).
\]
Es función acotada.

Demostración: ejercicio, usar la desigualdad del triángulo.

\section*{Foto 8}

\textit{Afirmación:} $D$ es una métrica en el espacio $\mathcal B(A,X)$.

\begin{enumerate}[label=\arabic*)]
    \item Positividad. Si $f,g\in\mathcal B(A,X)$, $f\neq g$, $\therefore \exists x_0\in A$ tal que $f(x_0)\neq g(x_0)$.
\end{enumerate}

\section*{Foto 9}

\[
D(f,g)=\sup\{d(f(x),g(x))\mid x\in A\}
\]
\[
\ge d(f(x_0),g(x_0))>0.
\]

Si $f=g$ entonces
\[
D(f,f)=\sup\{d(f(x),f(x))\mid x\in A\}=0.
\]

\begin{enumerate}[label=\arabic*)]
    \setcounter{enumi}{1}
    \item Simetría.
\end{enumerate}

\[
d(f(x),g(x))=d(g(x),f(x)),\ \forall x\in A.
\]
\[
\therefore D(f,g)=D(g,f).
\]

\begin{enumerate}[label=\arabic*)]
    \setcounter{enumi}{2}
    \item Desigualdad del triángulo. Sean $f,g,h\in\mathcal B(A,X)$.
\end{enumerate}

\[
D(f,g)\le D(f,h)+D(h,g).
\]

\section*{Foto 10}

\[
d(f(x),g(x))\le d(f(x),h(x))+d(h(x),g(x))
\]
por la desigualdad del triángulo en $(X,d)$.

\[
D(f,g)=\sup\{d(f(x),g(x))\mid x\in A\}
\]
\[
\le \sup\{d(f(x),h(x))\mid x\in A\}+\sup\{d(h(x),g(x))\mid x\in A\}.
\]
La desigualdad (?) se da por las mismas razones que en $C[a,b]$.

\section*{Foto 11}

\textit{Ejemplo.} Sea $A=\{1,2\}$, $X=\mathbb R$.
\[
\mathcal B(A,X)=\{f:A\to X\mid f \text{ es acotada}\}.
\]
Acotado en el sentido de hace un momento.

\textit{Afirmación:} $\mathcal B(A,X)$ = todas las funciones de $A$ en $X$.
Para cualquier $f:A\to\mathbb R$, es acotada porque
\[
f(A)=\{f(1),f(2)\}\subset \mathbb R.
\]
Hay que probar que $f$ es función acotada según \dots

\section*{Foto 12}

Según la definición de función acotada toma valores en un espacio métrico acotado.

\[
\exists g\in \overline X \text{ tal que } \gamma(a)=d(f(a),g).
\]

\section*{Foto 13}

Retomamos el ejemplo:
\[
A=\{1,2\},\qquad X=\mathbb R,
\]
\[
d \text{ es la métrica usual en } \mathbb R.
\]
\[
\mathcal B(A,X)=\{f:A\to\mathbb R\}=\{(f(1),f(2))\mid f:A\to\mathbb R\}=(\mathbb R)^2.
\]

\[
\text{¿Cómo se ve la métrica }D?
\]
\[
D(f,g)=\sup\{d(f(x),g(x))\mid x\in\{1,2\}\}.
\]

\section*{Foto 14}

\[
D(f,g)=\sup\{|f(1)-g(1)|,\ |f(2)-g(2)|\}.
\]
\[
=\max\{|f(1)-g(1)|,\ |f(2)-g(2)|\}.
\]

\[
D\bigl((f(1),f(2)),(g(1),g(2))\bigr).
\]
La métrica usual en $\mathbb R^2$:
\[
\bigl\|(f(1),f(2))-(g(1),g(2))\bigr\|_{\mathbb R^2}
\]
\[
=\sqrt{(f(1)-g(1))^2+(f(2)-g(2))^2}.
\]

\section*{Foto 15}

\textit{Definición.} Sea $(X,d)$ espacio métrico. Sean $p,q\in X$. Sea $r\in\mathbb R$, $r>0$.

\begin{enumerate}[label=\alph*)]
    \item Sea $p\in X$. Una vecindad o un entorno de $p$, de radio $r>0$, es el siguiente conjunto:
    \[
    N_r(p)=\{q\in X\mid d(q,p)<r\}.
    \]
    \item Sea $E\subset X$. Sea $p\in X$. Decimos que $p$ es punto límite de $E$.
\end{enumerate}

\section*{Foto 16}

Toda vecindad de $p$ contiene elementos de $E$ que no son $p$.
Es decir,
\[
\forall r>0,\quad N_r(p)\cap(E\setminus\{p\})\neq \varnothing.
\]

\textit{Ejemplo.} $X=\mathbb R$, $E=(0,1]\cup\{2\}$.
El $0$ y $1$ son puntos límite de $E$.

\section*{Foto 17}

\[
N_r(0)=\{q\in\mathbb R\mid |q|<r\}=(-r,r).
\]
\[
N_r\!\left(\tfrac12\right)=\{q\in\mathbb R\mid |q-\tfrac12|<r\}=
\left(\tfrac12-r,\tfrac12+r\right).
\]

\section*{Foto 18}

\[
\text{P.D. } \exists r>0\quad N_r(2)\cap (E\setminus\{2\})=\varnothing.
\]

\[
N_r(2)\cap E=\varnothing.
\]

\section*{Foto 19}

\textit{Continuamos con las definiciones.}

\begin{enumerate}[label=\alph*)]
    \item Si $p\in E$ y $p$ no es punto límite de $E$, decimos que $p$ es punto aislado de $E$.
\end{enumerate}

\textit{Ejemplo.} $X=\mathbb R$, $E=(0,1]\cup\{2\}$.
El $2$ es punto aislado de $E$.
El $2$ no es punto límite de $(0,1]$.

\section*{Fotos 20--23}

Continuamos con las definiciones:

\begin{enumerate}[label=\alph*)]
    \item $E\subset X$ es cerrado si contiene a todos sus puntos límite.
    \item Sea $p\in X$. Decimos que $p$ es punto interior de $E$ si $\exists r>0$ tal que $N_r(p)\subset E$.
    \item $E$ es conjunto abierto si todos sus elementos son puntos interiores.
    \item $E^c=\{x\in X\mid x\notin E\}$.
    \item $E\subset X$ se dice que es un conjunto perfecto si cumple con:
    \begin{enumerate}[label=\arabic*)]
        \item es cerrado;
        \item todos sus elementos son puntos límite de $E$.
    \end{enumerate}
    \item $E\subset X$ es conjunto acotado si $\exists M>0$ y $\exists q\in X$ tal que
    \[
    d(p,q)<M\qquad \forall p\in E.
    \]
    \item $E$ es denso en $X$ si todo $p\in X$ es punto límite de $E$.
\end{enumerate}

\textit{Ejemplo.} $X=\mathbb R$, $E=(0,1]\cup\{2\}$.
El $2$ es punto aislado de $E$.
El $2$ no es punto límite de $(0,1]$.

\end{document}
